\section{Enumerating the Spectrum}
\label{sec:mass_ledger}
\paragraph{Status of the ledger.} These are proposed geometric assignments. A map from the multiplicative mass template to the signed six-term ledger is still needed, together with a dimensionful mass anchor and a renormalization convention. The decay formalization does not certify these assignments as uniquely selected masses.


In the $E_8$-Persistence framework, mass is not an axiomatic parameter assigned to a point particle. It is the \textbf{Thermodynamic Latency} (entropic drag) incurred by an information packet as it executes its specific routing protocol across the discrete lattice. 

Because elementary fermions are persistent subsystems, their mass formulas must map perfectly to the universal Informational Energetics Impedance Ledger: 
\begin{equation}
    M_{particle} = \text{CAP} + \text{MAP} + \text{PRO} + \text{GOV} + \text{TOL} + \text{MAR}
\end{equation}

For fundamental, unconfined fermions (which have not yet undergone topological network closure), this ledger simplifies to three primary geometric factors:
\begin{equation}
    \frac{M_{\rm particle}}{m_{\rm ref}}=A_{\rm base}\,\Delta^n F_{\rm geom}
\end{equation}

\begin{itemize}
    \item \textbf{Base Anchor:} Reflects whether the protocol must anchor to the spacetime manifold. Boundary-coupled protocols (Leptons) have a base anchor of $1$. Color-coupled protocols (Quarks) must anchor their flux tubes to spacetime, incurring a base drag of $D=4$.
    \item \textbf{Resonant Tier ($\Delta^n$):} The temporal footprint of the standing wave across $n$ geometric dimensions. The fundamental lattice clock period is $\Delta=43$. Resonances exist as $0$D points ($\Delta^0=1$), $1$D lines ($\Delta^1$), $2$D areas ($\Delta^2$), and $3$D volumes ($\Delta^3$). 
\end{itemize}

\subsection{The Substrate Invariants (The Constants of Nature)}
To evaluate the geometric mass formulas, we utilize the strict, parameter-free invariants derived from the $E_8$ substrate in Papers I-III:
\begin{itemize}
    \item $\Delta = 43$: The fundamental temporal clock period (the Stark-Heegner resonance).
    \item $D = 4$: The macroscopic spacetime manifold dimension.
    \item $\sigma = 5$: The internal interaction symmetry rank.
    \item $\chi = 2$: The topological boundary (the Euler characteristic of the 2-sphere).
    \item $\nu = 16$: The active chiral channel capacity.
    \item $\alpha \approx 1/137.036$: The baseline geometric admittance of the vacuum.
    \item $\eta = 171/172$: The Coordinate Overhead ($\frac{D\Delta - 1}{D\Delta}$), representing the ratio of addressable payload to total manifold resolution.
\end{itemize}

\subsection{The Geometric Modifiers}
The modifiers encode exactly how each protocol navigates the finite capacity constraints of the substrate. These are not fitted parameters; they are strict architectural consequences of the substrate invariants ($\sigma=5, \chi=2, \nu=16$):

\begin{itemize}
    \item \textbf{Static Anchor ($+\sigma$):} A $0$D point lacks the spatial extent to rotate; it pays the static additive weight of the internal symmetry.
    \item \textbf{Rotational Drag ($\times\sigma$):} A $1$D linear standing wave physically sweeps through the internal interaction symmetry, multiplying its latency by $\sigma$.
    \item \textbf{Color Dilution ($\div (\sigma-\chi)$):} A $2$D membrane residing in the color bulk must distribute its flux isotropically across the 3 internal dimensions. The effective manifold latency is diluted.
    \item \textbf{Boundary Overflow ($+\chi/\sigma$):} A $2$D area that exceeds the spatial manifold capacity leaks excess flux into the boundary-to-symmetry ratio.
    \item \textbf{Capacity Ceiling ($+1/D$):} A $3$D standing wave completely fills the local node. Hitting the $4$D capacity ceiling forces reciprocal volumetric leakage.
\end{itemize}

By applying these strict geometric rules to the four base routing protocols, we exhaustively generate the exact masses of the 12 fundamental fermions.

% ---------------------------------------------------------
\CatchFileBetweenTags{\ElectronVal}{particles.tex}{ElectronVal}
\CatchFileBetweenTags{\ElectronEq}{particles.tex}{ElectronEq}
\CatchFileBetweenTags{\ElectronExperimentalValue}{particles.tex}{ElectronExperimentalValue}
\CatchFileBetweenTags{\ElectronSigma}{particles.tex}{ElectronSigma}

\subsection{Protocol 1: The Electron (Boundary-Coupled Singlets)}

This protocol bypasses the internal color bulk entirely, routing $100\%$ of its flux through the $\chi=2$ topological boundary. It has a base anchor of $1$.

\paragraph{Tier 0 (Point): The Electron.} The unitary ground state. It possesses no internal spatial dimensions and thus incurs no geometric drag.
\begin{align*}
    M_e &= \ElectronEq \cdot m_e \\
        &= \num[round-mode=figures,round-precision=9]{\ElectronVal} \text{ MeV}
\end{align*}

\CatchFileBetweenTags{\MuonVal}{particles.tex}{MuonVal}
\CatchFileBetweenTags{\MuonEq}{particles.tex}{MuonEq}
\CatchFileBetweenTags{\MuonExperimentalValue}{particles.tex}{MuonExperimentalValue}
\CatchFileBetweenTags{\MuonSigma}{particles.tex}{MuonSigma}

\paragraph{Tier 1 (Line): The Muon.} A linear string ($\Delta^1$). It sweeps through the internal symmetry ($\sigma\Delta$), minus the topological anchor ($\chi$) and the double-cover spin phase lag ($2\pi$). Because it is unconfined, it actively couples to the macroscopic vacuum. On a discrete network, this coupling is an infinite recursive geometric feedback loop. The total informational load transferred to the vacuum is the infinite geometric series of the admittance $\alpha$: $\sum \alpha^n = \frac{\alpha}{1-\alpha}$. 
\begin{align*}
    M_\mu &= \left( \MuonEq \right) \cdot m_e \\
          &\approx \num[round-mode=figures,round-precision=9]{\MuonVal} \text{ MeV} \quad (\text{Pull: } \num[round-mode=places,round-precision=2]{\MuonSigma}\sigma)
\end{align*}
\textit{Interpretive question (the anomalous magnetic moment).} A coefficient in a geometric mass expansion does not calculate the electromagnetic vertex form factor that defines $g-2$. A current matrix element and its corrections are required before this proposed feedback interpretation can be tested.

\textit{Structural Consequence 2 (The Decay Residual):} This extreme precision highlights the ultimate limit of the continuous approximation. Because the Muon protocol requires a continuous $2\pi$ phase rotation, it is subject to the exact same \textbf{Lattice Discretization Error} observed in the Hadrons ($\sim 5-30$ eV). Furthermore, because the Muon decays via the Weak interface, its routing state remains continuously coupled to the Neutrino channels. The $\sim 5$ eV residual perfectly matches the magnitude of the Neutrino thermal ground state, reflecting the zero-point vibration of the decay channels inherently coupled to the Muon state.

\CatchFileBetweenTags{\TauVal}{particles.tex}{TauVal}
\CatchFileBetweenTags{\TauEq}{particles.tex}{TauEq}
\CatchFileBetweenTags{\TauExperimentalValue}{particles.tex}{TauExperimentalValue}
\CatchFileBetweenTags{\TauSigma}{particles.tex}{TauSigma}

\paragraph{Tier 2 (Area): The Tau.} An areal membrane. It pays the topological boundary cost explicitly across the causal area ($\chi\Delta^2$). To persist, it must shed the latency of its lower-tier structures: the linear core latency ($\sigma\Delta$) and spacetime anchors ($D+\chi$).
\begin{align*}
    M_\tau &= \left( \TauEq \right) \cdot m_e \\
           &\approx \num[round-mode=figures,round-precision=9]{\TauVal} \text{ MeV} \quad (\text{Pull: } \num[round-mode=places,round-precision=2]{\TauSigma}\sigma)
\end{align*}

% ---------------------------------------------------------
\subsection{Protocol 2: The Up-Type Quarks (Color-Coupled, Base D)}

\CatchFileBetweenTags{\UpVal}{particles.tex}{UpVal}
\CatchFileBetweenTags{\UpEq}{particles.tex}{UpEq}
\CatchFileBetweenTags{\UpExperimentalValue}{particles.tex}{UpExperimentalValue}
\CatchFileBetweenTags{\UpSigma}{particles.tex}{UpSigma}

This fractional protocol distributes its flux through the internal color bulk. Its flux tubes must explicitly anchor to the spacetime manifold, incurring a base anchor of $D=4$.

\paragraph{Tier 0 (Point): The Up Quark.} The minimal fractional protocol anchored purely to the spacetime manifold.
\begin{align*}
    M_u &= \UpEq \cdot m_e \\
        &\approx \num[round-mode=figures,round-precision=9]{\UpVal} \text{ MeV} \quad (\text{Pull: } \num[round-mode=places,round-precision=2]{\UpSigma}\sigma)
\end{align*}

\CatchFileBetweenTags{\CharmVal}{particles.tex}{CharmVal}
\CatchFileBetweenTags{\CharmEq}{particles.tex}{CharmEq}
\CatchFileBetweenTags{\CharmExperimentalValue}{particles.tex}{CharmExperimentalValue}
\CatchFileBetweenTags{\CharmSigma}{particles.tex}{CharmSigma}

\paragraph{Tier 2 (Area): The Charm Quark.} An areal anchor ($D\Delta^2$). Because it is a $2$D membrane operating in the color bulk, its flux is strictly diluted by the 3 internal dimensions ($\sigma-\chi$).
\begin{align*}
    M_c &= \left( \CharmEq \right) \cdot m_e \\
        &\approx \num[round-mode=figures,round-precision=9]{\CharmVal} \text{ MeV} \quad (\text{Pull: } \num[round-mode=places,round-precision=2]{\CharmSigma}\sigma)
\end{align*}

\CatchFileBetweenTags{\TopVal}{particles.tex}{TopVal}
\CatchFileBetweenTags{\TopEq}{particles.tex}{TopEq}
\CatchFileBetweenTags{\TopExperimentalValue}{particles.tex}{TopExperimentalValue}
\CatchFileBetweenTags{\TopSigma}{particles.tex}{TopSigma}

\paragraph{Tier 3 (Volume): The Top Quark.} A volumetric anchor ($D\Delta^3$). The $3$D standing wave physically hits the absolute $4$D capacity ceiling of the local node, forcing reciprocal leakage ($+1/D$).
\begin{align*}
    M_t &= \TopEq \cdot m_e \\
        &\approx \num[round-mode=figures,round-precision=9]{\TopVal} \text{ MeV} \quad (\text{Pull: } \num[round-mode=places,round-precision=2]{\TopSigma}\sigma)
\end{align*}

% ---------------------------------------------------------
\subsection{Protocol 3: The Down-Type Quarks (Color-Coupled, Base D + $\sigma$)}

This fractional protocol shares the manifold anchor ($D$), but carries the additional geometric load of the internal interaction symmetry ($\sigma$).

\CatchFileBetweenTags{\DownVal}{particles.tex}{DownVal}
\CatchFileBetweenTags{\DownEq}{particles.tex}{DownEq}
\CatchFileBetweenTags{\DownExperimentalValue}{particles.tex}{DownExperimentalValue}
\CatchFileBetweenTags{\DownSigma}{particles.tex}{DownSigma}

\paragraph{Tier 0 (Point): The Down Quark.} The static spacetime anchor ($D$) plus the static internal symmetry weight ($\sigma$).
\begin{align*}
    M_d &= \left( \DownEq \right) \cdot m_e \\
        &\approx \num[round-mode=figures,round-precision=9]{\DownVal} \text{ MeV} \quad (\text{Pull: } \num[round-mode=places,round-precision=2]{\DownSigma}\sigma)
\end{align*}

\CatchFileBetweenTags{\StrangeVal}{particles.tex}{StrangeVal}
\CatchFileBetweenTags{\StrangeEq}{particles.tex}{StrangeEq}
\CatchFileBetweenTags{\StrangeExperimentalValue}{particles.tex}{StrangeExperimentalValue}
\CatchFileBetweenTags{\StrangeSigma}{particles.tex}{StrangeSigma}

\paragraph{Tier 1 (Line): The Strange Quark.} A linear anchor ($D\Delta$) augmented by an additive double-symmetry drag.
\begin{align*}
    M_s &= \left( \StrangeEq \right) \cdot m_e \\
        &\approx \num[round-mode=figures,round-precision=9]{\StrangeVal} \text{ MeV} \quad (\text{Pull: } \num[round-mode=places,round-precision=2]{\StrangeSigma}\sigma)
\end{align*}

\CatchFileBetweenTags{\BottomVal}{particles.tex}{BottomVal}
\CatchFileBetweenTags{\BottomEq}{particles.tex}{BottomEq}
\CatchFileBetweenTags{\BottomExperimentalValue}{particles.tex}{BottomExperimentalValue}
\CatchFileBetweenTags{\BottomSigma}{particles.tex}{BottomSigma}

\paragraph{Tier 2 (Area): The Bottom Quark.} An areal anchor ($D\Delta^2$). Unlike the Charm quark, the space dimension saturates, forcing the geometric latency to leak into the boundary-to-symmetry ratio ($\chi/\sigma$).
\begin{align*}
    M_b &= \BottomEq \cdot m_e \\
        &\approx \num[round-mode=figures,round-precision=9]{\BottomVal} \text{ MeV} \quad (\text{Pull: } \num[round-mode=places,round-precision=2]{\BottomSigma}\sigma)
\end{align*}

\subsection{The Observer Effect: QCD Scale Dilation ($\eta$)}
The $\sigma$-pulls for the unconfined leptons (Muon, Tau) demonstrate extreme geometric precision. Conversely, the fractional quarks exhibit consistent $1 - 3\sigma$ residuals. 

This is not a flaw in the geometric lattice; it is the \textbf{Observer Effect}. Experimental measurements are taken from the macroscopic $D=4$ spacetime manifold. For Color-Coupled protocols (Quarks), the information routes through the internal color bulk. To observe an internal state from the macroscopic manifold, the state's energy must be projected outward through the causal grid. 

As derived in Paper III, traversing the continuous manifold incurs the Coordinate Overhead ($\eta = 171/172$). The $E_8$-Persistence formulas derive the strict \textbf{Geometric Pole Mass}—the exact latency of the routing protocol isolated on the lattice node. The observed variance perfectly characterizes the QCD scale dilation ($1/\eta$) required to translate these absolute topological invariants into renormalized macroscopic parameters evaluated via the $\overline{\text{MS}}$ scheme.

% ---------------------------------------------------------
\subsection{Protocol 4: The Neutrinos (Unanchored Phonons)}

The Neutrino protocol bypasses the color bulk (Colorless) and possesses zero boundary flux footprint ($Q=0$). Because it lacks both the $D$ manifold anchor and the $\chi$ boundary anchor, it cannot form a localized standing wave. 

It must propagate as a delocalized lattice vibration (a Phonon). Consequently, neutrinos follow an \textbf{Inverse Resonance Hierarchy} ($M \propto \Delta^{-n}$). Their base geometric floor is the absolute thermal resolution limit of the lattice ($\nu_{base} = \frac{\pi \alpha}{\nu \Delta^3}$), generating masses in the sub-eV scale.

\CatchFileBetweenTags{\NeutrinovthreeVal}{particles.tex}{NeutrinovthreeVal}
\CatchFileBetweenTags{\NeutrinovthreeEq}{particles.tex}{NeutrinovthreeEq}

\paragraph{Tier 0 (Base Volume): $\nu_3$} The unanchored baseline thermal floor.
\begin{align*}
    M_{\nu 3} &= \left( \NeutrinovthreeEq \right) \cdot m_e \approx \NeutrinovthreeVal \text{ MeV}
\end{align*}

\CatchFileBetweenTags{\NeutrinovtwoVal}{particles.tex}{NeutrinovtwoVal}
\CatchFileBetweenTags{\NeutrinovtwoEq}{particles.tex}{NeutrinovtwoEq}

\paragraph{Tier 1 (Areal Inverse): $\nu_2$} 
\begin{align*}
    M_{\nu 2} &= \left( \NeutrinovtwoEq \right) \cdot m_e \approx \NeutrinovtwoVal \text{ MeV}
\end{align*}

\CatchFileBetweenTags{\NeutrinovoneVal}{particles.tex}{NeutrinovoneVal}
\CatchFileBetweenTags{\NeutrinovoneEq}{particles.tex}{NeutrinovoneEq}

\paragraph{Tier 2 (Linear Inverse): $\nu_1$} 
\begin{align*}
    M_{\nu 1} &= \left( \NeutrinovoneEq \right) \cdot m_e \approx \NeutrinovoneVal \text{ MeV}
\end{align*}

\paragraph{The Neutrino Mass Anomaly (Cosmology vs. Oscillations):}
Standard physics faces a severe tension: cosmological surveys constrain the sum of neutrino masses to near-zero ($\Sigma m_\nu < 120$ meV), while atmospheric oscillation experiments imply much heavier splittings ($\sim 50$ meV). 

In the $E_8$-Persistence framework, this is a contextual Observer Effect. Neutrinos are unanchored lattice phonons with an intrinsic Rest Mass dictated by the thermal resolution floor ($M_{rest} \approx 9.2$ meV). Cosmology measures this decoupled, free-streaming state. 

However, when a neutrino participates in an oscillation or weak decay, it is actively processed through the $W$-boson Protocol Layer. Passing through this dense interface subjects the phonon to a geometric refractive enhancement ($\eta_W = \sigma + 1 = 6$). The active latency of the interacting state is therefore $M_{weak} \approx M_{rest} \times 6 \approx 55$ meV. Squaring these two contextual limits gives $M_{weak}^2-M_{rest}^2=35M_{rest}^2$ between $2.967$ and $2.969\times10^{-3}$~eV$^2$, and $M_{weak}^2$ alone between $3.052$ and $3.054\times10^{-3}$~eV$^2$ (\texttt{NeutrinoContext.weak\_minus\_rest}, \texttt{weak\_squared\_excluded}). The measured atmospheric splitting is $\Delta m^2_{3\ell}=(2.534\pm0.025)\times10^{-3}$~eV$^2$ (NuFIT~6.0, normal ordering), so both readings overshoot by about $17\%$, more than $17\sigma$ (\texttt{weak\_minus\_rest\_excluded}). A common multiplier would have to be $5.55<k<5.57$ rather than $\sigma+1=6$ (\texttt{atmospheric\_factor}), and no common multiplier can fit both oscillation scales, since the tower keeps the gap ratio $1849$ against an observed ratio near $34$ (Paper~III, \texttt{FloorNeutrino.no\_moderate\_gap\_ratio}). The contextual reading therefore does not reproduce the oscillation data as stated; a derived propagation Hamiltonian is required.